\section{Преглед на съществуващите решения}
\label{sec:existing}

Предходната глава показа, че инжекциите са добре познат клас уязвимости, чрез който
се извършват престъпления и който води до санкции за администраторите, и че политиката се измества от наказване към изключване на
класа още при проектирането. Настоящата глава разглежда какви средства за това вече
съществуват и къде всяко от тях спира. Първо се разглеждат защитите по слоеве, после
академичната работа върху SQL и инжекциите извън SQL. Следва прегледът на
библиотеките за типово безопасно конструиране на заявки. Главата завършва със
сравнителна таблица и с формулировка на пропуска, който следващите глави адресират.

Прегледът се основава на документацията на проектите и на оригиналните публикации.
Когато в прочетените страници даден факт не е посочен, в таблиците е записано „н.д.“
(не е посочено). Авторът не е стартирал разглежданите библиотеки, затова твърденията
за тях са твърдения на документацията, а не измерване.

\subsection{Защити по слоеве}
\label{sec:existing-layers}

Основният практически източник е наръчникът на \textit{OWASP} за предотвратяване на
SQL инжекции~\cite{owasp_sqli_prevention}. Той поставя на първо място подготвените
изрази с параметризирани заявки, следвани от съхранените процедури и проверката по
списък на допустимите стойности. Екранирането е определено като „STRONGLY
DISCOURAGED“, а към защитите се добавят най-малките привилегии. За всяка защита по-долу
се разглеждат две неща: какво гарантира и как се проваля.

\textbf{Екраниране} е преобразуване на специалните знаци във входа така, че
интерпретаторът да ги третира като данни. Наръчникът настоява ръчното екраниране да не
замества свързването на параметри, защото правилата зависят от диалекта и от
конфигурацията на базата~\cite{owasp_sqli_prevention}. Гаранцията е локална: важи за
конкретния диалект и място, където екранирането е извършено, и се губи на всяко
пропуснато място. Екранирането остава практика там, където параметри не могат да се
ползват: ръководството на \textit{Neo4j} препоръчва за идентификаторите в обратни
апострофи тези знаци да се удвояват~\cite{neo4j_cypher_injection_kb}.

При \textbf{параметризираните заявки} текстът и стойностите се предават поотделно:
заявката съдържа маркери, а драйверът свързва стойностите към тях. Стойността не може
да промени структурата, защото тя вече е разчетена. Това е защитата, препоръчана от
\textit{OWASP}, \textit{CISA} и \textit{FBI}~\cite{owasp_sqli_prevention,cisa_sqli_alert}.
Същият принцип стои зад \textit{Neo4j}, където планът на заявката е компилиран преди
подставянето на стойностите~\cite{neo4j_cypher_injection_kb}. В \textit{Cassandra}
заявка с поне един маркер трябва да бъде подготвена; документацията на \textit{CQL}
описва това като оптимизация и не го свързва със сигурността~\cite{cassandra_cql_prepared}.
Границата е посочена в наръчника на \textit{OWASP}: маркерите заместват стойности, а
не имена на таблици и колони и не ключови думи~\cite{owasp_sqli_prevention}.
Документацията на \textit{EF Core} казва същото~\cite{ef_core_sql}. Защитата не се
прилага автоматично: ако заявката е сглобена с конкатенация, параметри няма. Тя не
отговаря и на въпроса дали потребителят има право да види даден запис.

\textbf{Съхранените процедури} пренасят заявките в базата. Наръчникът предупреждава,
че свързването на аргументите защитава извикването, но не и динамичния SQL, който
процедурата сглобява вътре в себе си~\cite{owasp_sqli_prevention}. Халфонд, Виегас и
Орсо описват същото като разпространено заблуждение~\cite{halfond2006classification}.
Гаранцията е условна: зависи от начина, по който е написана процедурата.

Проверката по \textbf{списък на допустимите стойности} (\textit{allow-list}) се
препоръчва за структурни избори, които не могат да се параметризират. Пример е
посоката на сортиране: входът се съпоставя с литералите \code{ASC} или \code{DESC},
притежавани от приложението. Проверката допълва свързването на параметри, но не го
замества~\cite{owasp_sqli_prevention}. Гаранцията е силна само когато множеството може
да се изброи и не помага при свободен текст.

\textbf{Най-малките привилегии} ограничават какво може да направи вече променена
заявка. \textit{OWASP} предупреждава приложението да не получава права на
администратор на базата~\cite{owasp_sqli_prevention}, а \textit{Neo4j} препоръчва
контрол на достъпа по роли~\cite{neo4j_cypher_injection_kb}. Мярката не предотвратява
инжекцията, а намалява последиците ѝ.

\textbf{Защитната стена за уеб приложения} (\textit{web application firewall},
\textit{WAF}) проверява HTTP трафика пред приложението. \textit{ModSecurity} е
отворен механизъм за такива стени за \textit{Apache}, \textit{IIS} и
\textit{Nginx}~\cite{modsecurity_repo}, а \textit{OWASP Core Rule Set} (\textit{CRS})
е наборът правила за него. Всяко съвпаднало правило добавя точки към резултата на
заявката, която се отхвърля при праг, по подразбиране 5~\cite{crs_anomaly}.
По-високите от четирите нива на параноя откриват повече атаки, но дават повече
фалшиви тревоги, които налагат изключения на правила и могат да доведат до
изключване на \textit{CRS}~\cite{crs_paranoia}. Стената вижда HTTP заявката, а не
заявката към базата, която приложението по-късно ще сглоби. Гаранцията ѝ е
вероятностна и не променя причината за уязвимостта.

Защитите се разделят на две групи. Екранирането, параметрите и списъкът на допустимите
стойности действат върху самата заявка и дават силна гаранция, но само за онази част,
за която са приложими, и само когато програмистът ги приложи. Привилегиите и
защитната стена действат извън заявката и ограничават вредата или вероятността.
Общото между всички е, че се прилагат от човек на всяко място поотделно.

\subsection{Академична работа по откриване и предотвратяване}
\label{sec:existing-academic}

Су и Васерман определят инжекцията формално: за да успее атака, входът, който попада
в заявката, трябва да променя нейната предвидена синтактична структура. Според тях
това е първото формално определение на атаките чрез инжектиране на команди в уеб
приложенията. Заедно с него авторите дават алгоритъм за предотвратяването им, основан
на контекстно-свободни граматики. Реализацията им \textit{SQLCHECK}, оценена върху
реални приложения със системно събрани реални атаки, по думите им не е дала нито
фалшиво положителни, нито фалшиво отрицателни резултати~\cite{su2006essence}. Определението
се използва и в настоящата работа (глава~\ref{sec:theory}).

Халфонд, Виегас и Орсо правят обзор на известните видове SQL инжекции (сред тях
тавтологии, обединени и натрупани заявки, атаки чрез съхранени процедури,
алтернативни кодировки) и оценяват съществуващите техники спрямо тях. Основното им
наблюдение е, че много предложени решения адресират само подмножество от
проблема~\cite{halfond2006classification}. Техниката \textit{AMNESIA} на Халфонд и
Орсо строи чрез статичен анализ модел на законните заявки на приложението и проверява
всяка генерирана заявка спрямо него по време на изпълнение. Върху седем приложения
тя е спряла всички опитани атаки без фалшиво положителни резултати~\cite{halfond2005amnesia}.
\textit{CANDID} на Бандхакави и съавтори е, според резюмето, схема за
автоматично преобразуване на уеб приложения, така че по твърдение на авторите те да
са безопасни срещу SQL инжекции. Методът е посочен в заглавието като динамична
оценка на кандидати~\cite{bandhakavi2007candid}. \textit{SQLrand} на Бойд и
Керомитис рандомизира набора от инструкции на SQL, така че инжектираният текст да е
невалиден за парсера, а междинен прокси превежда рандомизирания език към стандартния
за \textit{MySQL}~\cite{boyd2004sqlrand}.

Тези работи имат общи черти. Всички са насочени към SQL. Действат върху вече
построена заявка: сравняват я с модел, с кандидат или с рандомизиран речник, а не
променят начина, по който се построява. Изискват допълнителен инструмент: анализ,
модифицирана среда или прокси. Гаранцията се налага по време на изпълнение върху вече
построената заявка, а не при построяването ѝ. Някои от работите дават и формално
обосновани алгоритми (така Су и Васерман претендират за алгоритъм, който е коректен и
пълен), но оценката им е върху подбрани приложения. Определението на Су и Васерман
остава общата основа, спрямо която се измерва успехът на всички тях.

\subsection{Инжекции извън SQL}
\label{sec:existing-nosql}

Литературата опровергава мнението, че хранилищата без SQL са имунизирани. Рон,
Шулман-Пелег и Бронщайн представят техники за инжектиране в NoSQL хранилища, сред
които инжектиране чрез масиви в \textit{PHP}, \textit{OR} инжектиране в
\textit{MongoDB} и инжектиране на произволен \textit{JavaScript}. В един пример с
оператора \code{$ne} условията за потребителско име и парола се изпълняват за всяка
стойност, различна от посочената~\cite{ron2015nosql}. Част от същите автори публикуват
статия в \textit{IEEE Security \& Privacy}, според резюмето на която само анализът на
кода е недостатъчен за предотвратяване на такива атаки и мерките трябва да обхващат
целия жизнен цикъл на софтуера~\cite{ron2016nosql}. Ръководството за тестване на уеб сигурността на
\textit{OWASP} (версия 4.2) обяснява, че атаките срещу NoSQL могат да се изпълняват в
процедурен език, затова въздействието им може да е по-голямо, и че филтрирането на
знаци, специални за HTML, не защитава JSON интерфейс. В \textit{MongoDB} интерфейсът
очаква BSON и предоставя безопасно сглобяване на заявки, но несериализираният JSON и
изразите на \textit{JavaScript} се допускат в алтернативни параметри, като най-често
използваният е операторът \code{$where}~\cite{owasp_wstg42_nosql}.

Документацията на \textit{Neo4j} назовава Cypher инжекцията „братовчед“ на SQL
инжекцията и посочва като основна защита подаването на входа като параметри.
Конкатенирането на вход в низ, подаван на \code{apoc.cypher.doIt}, остава уязвимо,
а етикетите, типовете на връзките и имената на свойства според статията не могат да
се подават като параметри и затова се препоръчва да се избягват или санитизират~\cite{neo4j_cypher_injection_kb}.
Текущото ръководство по Cypher уточнява, че параметри не могат да заместват директно
ключове на свойства, типове връзки и етикети (например \code{n:$param} е невалидно),
но етикетите и типовете могат да се задават динамично чрез израз~\cite{neo4j_cypher_parameters}.
Следователно в Cypher имената остават избираеми по време на изпълнение по документиран
път.

Документацията на \textit{Redis} приема, че сървърът е за доверени клиенти в доверена
среда и че недоверен достъп трябва да се опосредства от слой, проверяващ входа. За
инжекциите тя твърди, че протоколът няма понятие за екраниране на низове, че „injection
is impossible under normal circumstances using a normal client library“ и че протоколът
използва низове с префикс за дължина, безопасни за двоични данни. Същото важи за
скриптовете на \textit{Lua}, но тялото на скрипт не бива да се сглобява от недоверени
низове~\cite{redis_security}. Условието „нормална клиентска библиотека“ означава, че
гаранцията на протокола не е гаранция на приложението. За \textit{Cassandra}
документацията на CQL описва маркерите и подготвените изрази, но в прочетените
страници не е открито самостоятелно ръководство за инжекции в CQL~\cite{cassandra_cql_prepared}.

Всяко хранилище решава проблема със свои средства: параметри, подготвени изрази,
протокол с префикс за дължина. Стойностите са защитени по замисъл на протокола или
интерфейса, но точно там, където приложението избира структура или име,
документацията предупреждава за граница. Гаранциите са формулирани поотделно от
производителя на всяко хранилище. В прегледаните източници няма общо определение,
което да обхваща всички хранилища, а академичните работи от предходния раздел не
покриват нито едно от тях.

\subsection{Типово безопасно конструиране на заявки}
\label{sec:existing-typesafe}

Друг подход е заявките да не се пишат като низове, а да се изграждат от типове и
функции на езика. Така компилаторът проверява имената и типовете. Разглеждат се
\textit{LINQ} с \textit{EF Core} (C\#), \textit{jOOQ} (Java), \textit{Diesel} и
\textit{sqlx} (Rust), \textit{esqueleto} и \textit{opaleye} (Haskell) и
\textit{sqlpp11} (C++). За всяка се търси какво прави безопасно и дали остава път към
суров текст или към идентификатори, избрани по време на изпълнение.

Според документацията на \textit{EF Core} методите \code{FromSql} и
\code{FromSqlInterpolated} са безопасни, защото предават данните като параметър, а
\code{FromSqlRaw} може да е уязвим. Когато името на колона трябва да се вмъкне
динамично, се ползва \code{FromSqlRaw}, а санитизирането е отговорност на
програмиста~\cite{ef_core_sql}. В \textit{jOOQ} SQL обикновено се строи като типизирано дърво със
свързани стойности и според документацията по този път не е възможно да се въведат инжекции. Изключение
са методите за суров SQL, отбелязани с \code{org.jooq.PlainSQL} от версия 3.6~\cite{jooq_sqli},
чието използване документацията свързва с проверка чрез \textit{Checker Framework} или
\textit{Error Prone}~\cite{jooq_plain_sql}. \textit{Diesel} заявява стремеж да
елиминира неправилните взаимодействия с базата на етапа на компилацията~\cite{diesel_site}
и според документацията изисква схемата да е описана в кода~\cite{diesel_crate}. Суров
SQL се пише чрез \code{sql\_query}, а малки фрагменти чрез \code{sql}; страницата на
\code{sql\_query} не споменава инжекциите и предупреждава само за типовете~\cite{diesel_sql_query}.
Макросът \code{query!} на \textit{sqlx} проверява SQL низа при компилация срещу жива
база или запазен кеш~\cite{sqlx_readme}. Във версия 0.9.0 функцията \code{sqlx::query}
приема \code{impl SqlSafeStr}, динамично построен низ трябва изрично да се обвие в
\code{AssertSqlSafe}, а документацията предупреждава срещу \code{format!}~\cite{sqlx_query_docs}.

В \textit{esqueleto} има \code{rawSql} с текст и списък от стойности, а SQL функция
може да се въведе чрез \code{unsafeSqlFunction} или \code{unsafeSqlBinOp}~\cite{esqueleto_hackage}.
В \textit{opaleye} функциите за таблица и колона приемат обикновен низ:
\code{table :: String -> ...} и \code{tableField :: String -> ...}~\cite{opaleye_table},
а модулът съдържа и \code{unsafeCast :: String -> ...}~\cite{opaleye_main}. Това не
доказва, че недоверени имена могат да бъдат подадени, но показва, че типовата система
не го забранява. \textit{sqlpp11} е типово безопасен вграден език за SQL в C++, чиито
таблици се пораждат от DDL и чиито заявки се проверяват при компилация за грешки в
синтаксиса, типовете и имената~\cite{sqlpp11_repo}. За неподдържани изрази
документацията описва \code{sqlpp::verbatim}~\cite{sqlpp11_functions}. Същият автор
поддържа и \textit{sqlpp23}, чиято документация описва \code{verbatim_table}, който
добавя таблици, дори когато те не са известни при компилация, както и \code{verbatim} и
\code{verbatim_clause}~\cite{sqlpp23_custom,sqlpp23_repo}. Това е най-близката
прегледана библиотека на C++ до подхода на настоящата работа и затова е важно, че и в
нея път към суров текст е предвиден по дизайн. Преглед, изчерпателен за всички
библиотеки на C++, не е правен.

\begin{table}[!htbp]
\centering
\caption{Типово безопасни библиотеки: какво защитават и къде остава изход}
\label{tab:libs}
\footnotesize
\setlength{\tabcolsep}{3pt}
\begin{tabularx}{\textwidth}{L{1.9cm}L{2.6cm}L{3.2cm}L{3.5cm}Y}
\toprule
\textbf{Библиотека} & \textbf{Стойности} & \textbf{Идентификатори и структура} & \textbf{Път към суров текст} & \textbf{Проверка} \\
\midrule
EF Core (LINQ) & Параметри при \code{FromSql}~\cite{ef_core_sql} & Имена на колони не се параметризират, нужен е \code{FromSqlRaw}~\cite{ef_core_sql} & \code{FromSqlRaw}, \code{SqlQueryRaw}, \code{ExecuteSqlRaw} & Типове при компилация, превод в SQL при изпълнение \\
jOOQ & Свързани стойности в дървото~\cite{jooq_sqli} & Типизирано дърво~\cite{jooq_sqli} & Методи с \code{@PlainSQL}~\cite{jooq_sqli,jooq_plain_sql} & Компилация (Java), по желание допълнителна проверка \\
Diesel & Свързани стойности~\cite{diesel_debug_query} & Схемата е описана в кода~\cite{diesel_crate} & \code{sql\_query}, \code{sql} & Компилация \\
sqlx & Свързани параметри~\cite{sqlx_query_docs} & \code{query!} проверява литерален низ срещу база~\cite{sqlx_readme} & \code{AssertSqlSafe} върху динамичен низ~\cite{sqlx_query_docs} & Компилация срещу жива база или кеш \\
esqueleto & \code{rawSql} приема стойности, за DSL н.д.~\cite{esqueleto_hackage} & н.д. & \code{rawSql}, \code{unsafeSqlFunction}, \code{unsafeSqlBinOp}~\cite{esqueleto_hackage} & Компилация (Haskell) \\
opaleye & н.д. & Имена като \code{String}~\cite{opaleye_table} & \code{unsafeCast} с низ~\cite{opaleye_main} & Компилация (Haskell) \\
sqlpp11 / sqlpp23 & В sqlpp11 подготвени изрази с \code{parameter()}~\cite{sqlpp11_prepared} & В sqlpp11 таблици от DDL~\cite{sqlpp11_repo}, в sqlpp23 е възможен \code{verbatim_table} & \code{verbatim}, \code{verbatim_table}, \code{verbatim_clause}~\cite{sqlpp11_functions,sqlpp23_custom} & Компилация (C++) \\
\bottomrule
\end{tabularx}
\end{table}

Таблица~\ref{tab:libs} показва повтарящ се модел. Защитата на стойностите е нормален
случай там, където документацията я посочва. Идентификаторите и структурата са
защитени дотолкова, доколкото програмистът остава в типизирания интерфейс. Във всяка
прегледана библиотека обаче е документиран път към суров текст. Обозначен е различно:
наставката \code{Raw} в \textit{EF Core}, анотацията \code{@PlainSQL} с възможност за
проверка в \textit{jOOQ}, явното \code{AssertSqlSafe} в \textit{sqlx}. Тези
обозначения помагат при преглед, но налагането им е конвенция или външен инструмент, а
не свойство на типовата система. В \textit{opaleye} именуването чрез низове е част от
обикновения интерфейс. Таблицата не казва, че библиотеките са небезопасни, а че
гаранцията им зависи от това, че изходът няма да бъде използван.

\subsection{Сравнение на подходите}
\label{sec:existing-comparison}

Таблица~\ref{tab:comparison} съпоставя разгледаните подходи по критерии, които
следват от определението на Су и Васерман~\cite{su2006essence}: кои части на заявката
се защитават, покриват ли се хранилища извън SQL, кога се налага гаранцията и има ли
изход. Оценките са на автора, въз основа на цитираните по-горе източници.
„Частично“ означава, че гаранцията важи при определени условия, посочени в
съответния раздел.

\begin{table}[!htbp]
\centering
\caption{Сравнение на подходите по критерии}
\label{tab:comparison}
\footnotesize
\setlength{\tabcolsep}{3pt}
\begin{tabularx}{\textwidth}{L{2.4cm}L{2.0cm}L{2.5cm}L{2.0cm}L{1.9cm}L{1.9cm}Y}
\toprule
\textbf{Подход} & \textbf{Стойности} & \textbf{Иденти-\newline фикатори} & \textbf{Структура} & \textbf{Не-SQL хранилища} & \textbf{Кога се налага} & \textbf{Изход} \\
\midrule
Екраниране & частично & частично & не & зависи от диалекта & изпълнение & всяко пропуснато място \\
Подготвени изрази & да & не & не & да, където има маркери & изпълнение & сглобяване на низ \\
Съхранени процедури & при параметри вътре & не & не & н.д. & изпълнение & динамичен SQL вътре \\
Списък на допустимите & да, за изброими & да, за изброими & да, за изброими & да & изпълнение & свободен текст \\
Най-малки привилегии & ограничава вредата & ограничава вредата & ограничава вредата & да, като принцип & внедряване & не е приложимо \\
Защитна стена & по вероятност & по вероятност & по вероятност & според правилата & внедряване & пропуснати атаки \\
Откриване по структура (SQLCHECK, AMNESIA, CANDID, SQLrand) & да & да & да & не, само SQL & изпълнение, с инструмент & само SQL в оценките \\
Типови езици за заявки & да & зависи от библиотеката & да, без суров текст & основно SQL & компилация & да, във всички прегледани \\
\bottomrule
\end{tabularx}
\end{table}

Таблица~\ref{tab:comparison} позволява три наблюдения. Първо, никой подход не
съчетава защита на стойностите, идентификаторите и структурата, покритие на хранилища
извън SQL и липса на изход. Параметрите покриват стойностите, но по природа не покриват
имената. Откриването по структура покрива всички части на SQL заявката, но е насочено
само към SQL и действа върху вече построена заявка. Типовите езици за заявки са
единствените, които налагат гаранцията при компилация, но във всички прегледани
случаи оставят път към суров текст. Второ, подходите, които действат при изпълнение
или внедряване, откриват или блокират проблема, след като кодът вече съществува, а
подходите при компилация могат да предотвратят самото изграждане на небезопасната
заявка, но гаранцията им е толкова силна, колкото е затворен изходът. Трето, ако
изходът е част от гаранцията, нито един прегледан подход не дава гаранция за целия
клас, а само за частта, която съответства на неговия механизъм. Тази разлика между
защита на отделните случаи и изключване на класа е същността на аргумента на
\textit{CISA} и \textit{FBI} (раздел~\ref{sec:cybercrime-prevention}).

\subsection{Пропускът}
\label{sec:existing-gap}

В рамките на прегледаната литература и документация могат да се направят три
наблюдения. Първо, работата по откриване и предотвратяване на инжекции се
концентрира върху SQL~\cite{su2006essence,halfond2006classification,halfond2005amnesia,bandhakavi2007candid,boyd2004sqlrand}.
Работата върху NoSQL е по-малко и описва техники и мерки за смекчаване~\cite{ron2015nosql,ron2016nosql},
а документацията за Neo4j, Redis и Cassandra е формулирана отделно за всяко
хранилище. Второ, типовите езици за заявки защитават стойностите, но в прегледаните
библиотеки запазват начин да се подаде идентификатор или текст на заявката по време на
изпълнение. Той е обозначен по различен начин, но е налице. Трето, \textit{SQLCHECK} и
\textit{AMNESIA} са измерени върху реални приложения, но само за SQL~\cite{su2006essence,halfond2005amnesia}.
В прегледаните източници не е открита оценка, която прилага гаранция, наложена при
компилация, към няколко разнородни хранилища и я сравнява с измерена уязвима базова
реализация.

Пропускът е следователно в съчетанието на три условия, а не във всяко от тях
поотделно: гаранция, наложена при компилация, обхващаща няколко разнородни хранилища
и оценена спрямо измерена уязвима базова реализация. Формулировката е ограничена до
прегледаната литература и не е твърдение за абсолютна новост. Следващите глави
адресират този пропуск: глава~\ref{sec:theory} въвежда понятията, а останалите
разглеждат подхода, неговата реализация и оценката му.
